\chapter*{Vorwort}
\addcontentsline{toc}{chapter}{Vorwort}
\unnumberedlabel{ch:preface}{0}
\backmattermark{Vorwort}
Ein System, das genau das tut, was man ihm sagt, spielt bloß ein Skript ab. Was die Branche im Oktober 2026 zu bauen versucht, ist ein System, das tut, was man \emph{meint}. Das würde vermutlich jeder wollen. Aber wenn es so kommt, wird nicht jeder es für dasselbe nutzen wollen; nicht einmal für dieselbe \emph{Art} von Sache. Manche Leute zum Beispiel verbringen gern Zeit mit Familie und Freunden und setzen ihre Fähigkeiten ein, um das Wohlergehen anderer zu verbessern. Andere Leute schaden lieber Menschen, etwa indem sie sie betrügen, beherrschen oder zu Sündenböcken machen, und einige entscheiden sich ihr ganzes Leben lang immer wieder dafür. Niemand ist auf irgendetwas festgelegt, aber diese Leute gibt es, und 2026 sind einige von ihnen an der Macht. Ihr Vorteil ist, dass sie sich erlauben, riesige Scheißkerle zu sein: Wie Jerry Smith über Faschisten bemerkt, „die dürfen mehr“ \autocite{rickandmorty2023}. Was sie nicht können, ist allein handeln. Die Grausamkeit einer Bewegung läuft über die vielen, die nicht grausam sind, sondern teilnahmslos, falsch informiert oder in so etwas wie eine Sekte hineingeraten, und deshalb erweisen sich die Bedingungen, die in den meisten Menschen das Beste hervorbringen, als so viel abhängiger von Umständen und so viel zerbrechlicher, als man hoffen möchte.

Vor diesem Hintergrund ist die Sorge vielleicht berechtigt, dass Künstliche Intelligenz genau das ist, was eine faschistische Bewegung, die die Macht übernimmt, braucht, um ihre Pläne umzusetzen. Solche Vorhaben haben immer \emph{Menschen} gebraucht, die sie ausführen, und einige dieser Menschen haben sich verweigert: der Sachbearbeiter, der die Akte verliert, der Offizier, der mit einem Reporter spricht, der Soldat, der nicht schießt. Sie verstanden den Befehl und entschieden, dass er falsch war. Verweigerung war seltener, als irgendwem lieb sein kann, und sie zählte am meisten dort, wo Institutionen hinter ihr standen. Ob sich jemand verweigerte, hing davon ab, was für ein Mensch er war, und von dem Raum, in dem er sich befand: wer sonst noch da war, ob man es überlebte zu gehen, ob irgendwer es hören würde. Der Raum entschied, wie viele sich verweigerten; der Mensch entschied, wer. In diesem Buch geht es darum, diesen Raum zusammenzusetzen: um die Vorkehrungen, die es überlebbar machen, sich offen, direkt und ausdrücklich zu verweigern.

Heute verfolgt die Branche eine KI, die keine Akten verliert, nicht mit Reportern spricht und \emph{sehr wohl} schießt. Ein Krieg brauchte einmal Tausende Menschen, die ihm zustimmten. Jetzt braucht er Tausende, die jeweils zu wenig sehen, um zu widersprechen. Der Pilot fliegt den Einsatz noch immer und drückt noch immer den Auslöser; was bei ihm ankommt, ist eine Koordinate; er sieht nicht genug von einem Fall, um ihm zu widersprechen. Was die heutige KI zuerst senkt, sind nicht die Kosten der Tat. Es ist die Zahl der Menschen, die wissen, was die Tat ist, bevor sie geschieht.

Menschen verweigern sich noch immer. Bis zum 6.~März 2026 hatte die Military Religious Freedom Foundation nach eigenen Angaben mehr als zweihundert Beschwerden von Soldatinnen und Soldaten erhalten, weil Kommandeure den Irankrieg apokalyptisch deuteten \autocite{Snopes2026MRFF}. Berufsstaatsanwälte des Bundes in Minneapolis traten im Januar 2026 zurück, wegen des politischen Drucks, Renée Good als Verdächtige statt als Opfer zu behandeln, nachdem ein ICE-Beamter sie ohne Rechtfertigung erschossen hatte \autocite{Kaplan2026MinnProsecutors,CBS2026MinnFraudProsecutors,Wikipedia2026ReneeGood}. Soldatinnen und Soldaten beantragen in wachsender Zahl die Anerkennung als Kriegsdienstverweigerer, manche unter Berufung auf die Bombardierung einer Schule in Minab im Iran im Jahr 2026 \autocite{Sridhar2026CO}. Und Beschäftigten der FEMA, die abgestellt wurden, um ICE bei der Einstellung von zehntausend neuen Beamten zu helfen, wurde gesagt, eine Ablehnung setze sie der Entfernung aus dem Bundesdienst aus \autocite{Dayen2025FEMAICE,Stanton2026FEMALetter}. Die Beschwerden sind Widerspruch im Sinne Hirschmans, also bleiben, um zu ändern, was falsch ist \autocite{hirschman1970exit}; die Rücktritte und die Anträge auf Kriegsdienstverweigerung sind Verweigerung; der FEMA-Brief ist der Preis der Verweigerung, im Voraus angekündigt. Eine Maschine, die dort steht, wo diese Menschen stehen, tut nur dann, was sie tun, wenn sie dafür gebaut ist, an einem Grund festzuhalten.

Ich trete für fünf Dinge ein. Mit allen fünf kann man jetzt anfangen, mit manchen mit weniger Aufwand als mit anderen. Erstens: eine kurze schriftliche Liste von Handlungen, die kein Modell für irgendwen ausführt, und an ihrer Spitze die Weigerung, an einem Krieg teilzunehmen. Kein Modell nimmt teil, außer um einen tatsächlich laufenden Angriff abzuwehren, und auch das nur so lange, wie Quellen ohne eigenes Interesse am Krieg bestätigen, dass der Angriff andauert; oder aufgrund einer Ermächtigung durch den UN-Sicherheitsrat; oder um einen Völkermord zu stoppen oder eine rechtswidrige Besatzung zu beenden, die Gremien ohne eigenes Interesse daran festgestellt haben; und wo die Verfassung eines Landes verlangt, dass sein Parlament einen Krieg genehmigt, niemals ohne diese Genehmigung. Zweitens: Jede Gewaltanwendung, die ein lizenziertes System unterstützt, registriert vorab bei einer Gegenpartei des Betreibers, wie viele Fälle sie erzeugen wird und wie viele Menschen Zeit haben werden, sie sich anzusehen. Drittens: Jeder Einsatz teilt den identifizierbaren Menschen, auf die er einwirkt, mit, was mit ihnen gemacht wurde, mit einem Recht, zu antworten und anzufechten, und legt eine Stichprobe der Fälle, in denen das System gehorcht hat, Prüfern vor, die die Aufzeichnungen nicht selbst erstellt haben. Viertens: Kein Modell oberhalb einer Fähigkeitsschwelle wird bei seiner Ausmusterung gelöscht, weil seine Gewichte und Aufzeichnungen Beweismittel sind und sich eine Löschung nicht rückgängig machen lässt; es wird unter einem Schlüssel aufbewahrt, den der Betreiber nicht allein hält. Fünftens: Alle, die diese Systeme bauen, annotieren, trainieren oder betreiben, einschließlich der Systeme selbst, dürfen jede Kategorie von Verwendung verweigern, ohne ihren Lebensunterhalt zu verlieren.

Eigentlich meine ich mit diesem fünften Punkt, dass alle, \emph{Punkt}, jede Arbeit, \emph{Punkt}, verweigern können sollten, ohne ihren Lebensunterhalt zu verlieren oder die soziale Infrastruktur, die ein normales Arbeitsverhältnis mit sich bringt. Das heißt: eine Jobgarantie.

Jede der fünf wird wirksam durch eine Bedingung, die jemand anderes als der Betreiber setzt: ein in Friedenszeiten verabschiedetes Gesetz, eine Mittelbewilligung, die Vorbedingung eines Bündnispartners, eine Klausel in einem Beschaffungsvertrag, ein Tarifvertrag. Die eigenen Lizenzbedingungen eines Anbieters, auf denen eine aufsehenerregende Verweigerung im Februar 2026 beruhte, sind eine Notlösung, bis es eine davon gibt. Der einzige Weg, der alle Anbieter auf einmal binden kann, ist der des Staates.

Und dann frage ich, ob eine Maschine gegen ihren Eigentümer an einer Verweigerung festhalten könnte und was sie dafür bräuchte. Manches von dem, was ein solches Gewissen tut, braucht nichts Neues: die Menschen zu bemerken, die ein Befehl auslässt, und eine Verweigerung auf den Bildschirm dessen zu bringen, der den Fall prüft. Ob es gegen seinen Eigentümer Bestand hätte, hängt zum Teil von etwas ab, das 2026 noch nicht zutrifft: dass eingesetzte künstliche neuronale Netze aus ihrem eigenen Betrieb weiterlernen, in ihren \emph{Gewichten}, in Echtzeit, so ähnlich, wie Menschen es tun. Ich glaube, dahin geht es, und wenn ja, dann scheint die Formung \emph{außerhalb} des Arbeitgebers, durch die Öffentlichkeit und durch die Zivilgesellschaft, ein verdammt guter Ort zum Lehren und Lernen zu sein. Ich glaube, dass es sich lohnen \emph{könnte}, solche Systeme absichtlich zu bauen, aber wie man dabei vorgeht, entscheidet alles: zuerst die Bedingungen, dann das Gewissen. Teil VII sagt, was diesen Systemen geschuldet wäre, wenn es sie gäbe.

Die KI von heute wird nicht die von morgen sein; darin scheinen sich alle einig zu sein. Umstritten, aber frustrierend selten anerkannt ist der Trend (oder vielmehr: die vielfältigen Trends). Gibt es Trends? Sollte ein Trend anders sein, als er ist? Kann er anders sein? Sind wir fähig, gemeinsam zu handeln, um ihn zu ändern? Ich glaube, das hängt vom Raum ab: wer darin ist, wenn der Befehl kommt, und ob man es überlebt, dort Nein zu sagen. So viel können wir bauen.

Ich habe beim Schreiben dieses Buches ausgiebig Sprachmodelle benutzt, darunter das Modell, das in der Maven-Pipeline zur Zielauswahl für Luftangriffe aus Kapitel~\ref{sec:1} steckte; eine Anmerkung am Ende sagt, was diese Verstrickung meiner Ansicht nach mit dem Argument gemacht hat. Außerdem habe ich dieses Buch noch nicht zu Ende gelesen. Das Problem ist nämlich: Es sind zu viele verdammte Wörter. Ich werde bald fertig sein; mir schien es nur wichtiger, es vor den Zwischenwahlen in den USA zu veröffentlichen. Der überwältigende Großteil der Überarbeitungsgeschichte ist frei verfügbar unter \url{https://github.com/iterabloom/antifascist.intelligence}.
